% !TeX encoding = UTF-8
% !TeX spellcheck = en_EN-EnglishUnitedKingdom
%% =======================================================
%%		P3C.TeX				v0.3
%%		SCVRIA SOFWARE		Ultimate UOC PEC Repository
%% =======================================================
%%	This file is part of the package P3C.TeX
%%
%%	Author's name: Scvria
%%	File: pxUML.code.tex
%% 	Version: v0.1 (02/2026)
%%	Description: Provides a UML class diagram implementation based on top of pgf-umlcd and TikZ packages
%%
%%	This file may be distributed and/or modified:
%%
%%	 1. under the LaTeX Project Public License and/or
%%	 2. under the GNU Public License.
%%
%%	 See the file doc/LICENSE for more details.
%%
%% =======================================================
%%	TODO:	- User level API
%%			- Add more features to the UML class diagram implementation
%% =======================================================

\documentclass[fake,cover-config=offset,default]{P3CTeX}
%\PECTeXconfig{cover-config={fullcover}}
%\documentclass{article}
%\usepackage[UML]{pxCORE}
\begin{document}

	\section{Proves bàsiques de pxPRP}


	% --- Prova 1: smoke test, creació i accés bàsic --------------------------------

	\subsection{Prova 1: Creació i accés bàsic}
	\verb*|\pxNewObject{PATATA}|\\
	\pxNewObject{PATATA}

	\verb*|\pxSetValue{PATATA}{k1}{v1,v2}|\\
	\pxSetValue{PATATA}{k1}{v1,v2}

	\verb*|\pxSetValue{PATATA}{k2}{vA,vB,vC}|\\
	\pxSetValue{PATATA}{k2}{vA,vB,vC}

	\verb*|\pxGET PATATA.k1[2]| $\to$ \pxGET PATATA.k1[2]\\

	\verb*|\pxPRINTF PATATA.k2[2]|$\to$ \pxPRINTF PATATA.k2[2]

	\bigskip

	\subsection{Prova 2: Múltiples objectes i claus aïllades}

	\pxNewObject{OBJAA}
	\pxNewObject{OBJBB}

	\pxSetValue{OBJAA}{key}{AA1,AA2}
	\pxSetValue{OBJBB}{key}{BB1,BB2}

	En pantalla: OBJAA.key[1] = \pxGET OBJAA.key[1],\quad
	OBJBB.key[1] = \pxGET OBJBB.key[1].

	\bigskip

	\subsection{Prova 3: Sobreescriptura amb \texttt{\textbackslash pxSet}}

	\pxNewObject{OVER}
	\pxSetValue{OVER}{state}{old}

	Abans de sobreescriure: OVER.state[1] = \pxGET OVER.state[1].

	\pxSetValue{OVER}{state}{newA,newB}

	Després de sobreescriure: OVER.state[1] = \pxGET OVER.state[1],\quad
	OVER.state[2] = \pxGET OVER.state[2].

	\bigskip

	\subsection{Prova 4: Acumulació de llistes amb \texttt{\textbackslash pxPutValue}}

	\pxNewObject{LST}

	\pxPutValue{LST}{items}{a}
	\pxPutValue{LST}{items}{b}
	\pxPutValue{LST}{items}{c}

	Resultat esperat: LST.items = (\pxGET LST.items[1], \pxGET LST.items[2], \pxGET LST.items[3]).

	\bigskip
	%
	\subsection{Prova 5: Esborrar claus i objectes}

	\pxNewObject{TEMP}
	\pxSetValue{TEMP}{k}{X,Y}

	Abans d'esborrar: TEMP.k[1] = \pxGET TEMP.k[1].

	\pxClearKey{TEMP}{k}

	Després d'esborrar la clau, podem tornar a definir-la:

	\pxSetValue{TEMP}{k}{Z}

	Ara: TEMP.k[1] = \pxGET TEMP.k[1].

	\pxClearObject{TEMP}

	\pxNewObject{TEMP}
	\pxSetValue{TEMP}{k}{P,P}
	Ara: TEMP.k[1] = \pxGET TEMP.k[1].
	\bigskip

	\subsection{Prova 6: Flags booleans}

	\pxNewObject{FLAGS}

	\pxSetValue{FLAGS}{featureA}{true}
	\pxSetValue{FLAGS}{featureB}{{false}}

	\pxIfTrueTF{FLAGS}{featureA}{ACTIVA}{INACTIVA}

	\par

	\pxIfTrueTF{FLAGS}{featureB}
	{featureB és ACTIVA.}
	{featureB és INACTIVA.}

	\bigskip

	\subsection{Prova 7 (mínima): Llistes i impressió completa}

	\pxNewObject{OBJL}
	\pxListSet{OBJL}{nums}{1,2,3}

	Tots els valors de OBJL.nums (esperat 1--3): \pxListAll{OBJL}{nums}[ -- ].

	\bigskip

	\pxNewObject{OBJS}
	\pxListPutRight{OBJS}{vals}{a}
	\pxListPutRight{OBJS}{vals}{b,c}

	Tots els valors de OBJS.vals (esperat a,b,c): \pxListAll{OBJS}{vals}.

	\bigskip

	\subsection{Prova 8 (mínima): Predicats d'existència i llargada}

	\pxIfObjectExistsTF{NOEXISTEIX}
	{OBJECTE NOEXISTEIX (no hauria de passar).}
	{OBJECTE NOEXISTEIX no és present, com esperat.}

	\pxNewObject{CHK}
	\pxSetValue{CHK}{nums}{1,2,3}

	\pxIfObjectExistsTF{CHK}
	{CHK existeix.}
	{CHK \emph{no} existeix (error).}

	\pxIfKeyExistsTF{CHK}{nums}
	{nums existeix a CHK.}
	{nums \emph{no} existeix a CHK (error).}

	Llargada esperada de CHK.nums: 3.\quad
	Llargada real: \pxListLen{CHK}{nums}.

	\bigskip

	\subsection{Prova 9: Accés escalar directe amb \texttt{\textbackslash pxGetScalar}}

	\pxNewObject{SCALAR}
	\pxSetValue{SCALAR}{flag}{true}
	\pxSetValue{SCALAR}{text}{Hola, món}

	Valor escalar de SCALAR.flag (esperat \texttt{true}): \pxGetScalar{SCALAR}{flag}.\\
	Valor escalar de SCALAR.text (esperat ``Hola, món''): \pxGetScalar{SCALAR}{text}.\\
	Valor escalar per a clau inexistent (esperat buit): [\pxGetScalar{SCALAR}{noexisteix}].

	\bigskip

	\subsection{Prova 10: Neteja de llistes amb \texttt{\textbackslash pxListClear}}

	\pxNewObject{LCLR}
	\pxListSet{LCLR}{vals}{a,b,c}

	Abans de netejar, llargada esperada 3, real: \pxListLen{LCLR}{vals}.\\
	\pxListClear{LCLR}{vals}

	Després de netejar, llargada esperada 0, real: \pxListLen{LCLR}{vals}.\\
	Es pot tornar a omplir:
	\pxListPutRight{LCLR}{vals}{x,y}
	; valors actuals (esperat x,y): \pxListAll{LCLR}{vals}.
	\section{Proves de la API pxUML}

	En aquesta secció es comprova l'ús de totes les ordres d'alt nivell del paquet.

	% Classe simple amb atributs i operacions modificats posteriorment
	\pxUMLClass{Person}
	{
		- name : String
	}
	[+ getName() : String]

	\addAttribute{Person}{- age : Integer}
	\addOperation{Person}{+ setName(newName : String) : void}

	% Classe base i derivada per provar la herència
	\pxUMLClass{BaseClass}
	{
		- id : Integer
	}

	\pxUMLClass{DerivedClass}
	{
		- extra : String
	}

	\addInheritance{DerivedClass}{BaseClass}

	% Classe utilitzada per instàncies
	\pxUMLClass{Service}
	{
		- name : String
	}
	[+ run() : void]

	\pxUMLDiagram{
		% dibuix complet
		\drawclass[text width=14em]{Person}[0,0]
		% dibuix simplificat
		\drawsimpleclass*[text width=10em]{BaseClass}[8,0]
		\drawsimpleclass[text width=10em]{DerivedClass}[8,-4]
		%	% instància d'una classe amb atributs i operacions
		\drawinstance*[14]{ServiceInstance}[0,-6]{Service}{
			- name : String
		}{
			+ run() : void
		}
	}[0.9\linewidth]
	% prova que \PXshowclass no produeix cap error en accedir al registre
	\PXshowclass{Person}
	\PXshowclass{Service}

	% inspecció bàsica: accessos de lectura no haurien de fallar
	Class Person attributes: [\pxUMLAttributes{Person}].\\
	Class Service operations: [\pxUMLOperations{Service}].\\
	Parents of DerivedClass: [\pxUMLParents{DerivedClass}].\\

	% cas negatiu: classe inexistent no ha de provocar error ni produir sortida
	\PXshowclass{NoExisteix}
	Class NoExisteix attributes (hauria de ser buit): [\pxUMLAttributes{NoExisteix}].

	\subsection{Diagrama UML}
	Disseny inicial del sistema, tal com s'indica a l'enunciat:

	\pxUMLClass[enumclass]{Role}
	{
		DEVELOPER,
		COMMUNITYMANAGER
	}
	\pxUMLClass[enumclass]{AICategory}
	{
		CHAT,
		CODEGENERATOR,
		POSTGENERATOR
	}

	\pxUMLClass{Employee}
	{
		- role : Role,
		- email : String
	}[+ monthlyConsumptionCost(month : DateTime) : real]

	\pxUMLClass{TokenConsumption}
	{
		- timeStamp : DateTime,
		- numTokens : UnlimitedNatural
	}

	\pxUMLClass{AIAgent}
	{
		- name : String,
		- description : String,
		- tokenPrice : Real,
		- aiCategory : AICategory
	}[+ process(prompt : String) : String, + calculateCost(numTokens : UnlimitedNatural):Real	]

	\pxUMLClass{Company}
	{
		- name : String,
		- address : String,
		- tel : String,
		- web : String
	}

	\pxUMLDiagram{
		\drawclass*[text width=15em]{Role}[2,0]
		\drawclass*[text width=14em]{AICategory}[16,0]
		\drawclass[text width=16em]{TokenConsumption}[9,-5]
		\drawclass[text width=16em]{AIAgent}[18,-5]
		\drawclass[text width=16em]{Employee}[0,-5]
		\drawclass[text width=16em]{Company}[0,-10]
		\association{Employee}{0..*}{ }{Company}{}{1}
		\association{Employee}{1}{ }{TokenConsumption}{}{0..*}
		\association{TokenConsumption}{0..*}{}{AIAgent}{}{1}
	}[1.1\linewidth]
	\subsection{Proves de la API pxUML}

	En aquesta secció es comprova l'ús de totes les ordres d'alt nivell del paquet.

	% Classe simple amb atributs i operacions modificats posteriorment
	\pxUMLClass{Person}
	{
		- name : String
	}
	[+ getName() : String]

	\addAttribute{Person}{- age : Integer}
	\addOperation{Person}{+ setName(newName : String) : void}

	% Classe base i derivada per provar la herència
	\pxUMLClass{BaseClass}
	{
		- id : Integer
	}

	\pxUMLClass{DerivedClass}
	{
		- extra : String
	}

	\addInheritance{DerivedClass}{BaseClass}

	% Classe utilitzada per instàncies
	\pxUMLClass{Service}
	{
		- name : String
	}
	[+ run() : void]

	\pxUMLDiagram{
		% dibuix complet
		\drawclass[text width=14em]{Person}[0,0]
		% dibuix simplificat
		\drawsimpleclass*[text width=10em]{BaseClass}[8,0]
		\drawsimpleclass[text width=10em]{DerivedClass}[8,-4]
		%	% instància d'una classe amb atributs i operacions
		\drawinstance*[14]{ServiceInstance}[0,-6]{Service}{
			- name : String
		}{
			+ run() : void
		}
	}[0.9\linewidth]
	% prova que \PXshowclass no produeix cap error en accedir al registre
	\PXshowclass{Person}
	\PXshowclass{Service}

	% inspecció bàsica: accessos de lectura no haurien de fallar
	Class Person attributes: [\pxUMLAttributes{Person}].\\
	Class Service operations: [\pxUMLOperations{Service}].\\
	Parents of DerivedClass: [\pxUMLParents{DerivedClass}].\\

	% cas negatiu: classe inexistent no ha de provocar error ni produir sortida
	\PXshowclass{NoExisteix}
	Class NoExisteix attributes (hauria de ser buit): [\pxUMLAttributes{NoExisteix}].

	\subsection{Diagrama UML}
	Disseny inicial del sistema, tal com s'indica a l'enunciat:

	\pxUMLClass[enumclass]{Role}
	{
		DEVELOPER,
		COMMUNITYMANAGER
	}
	\pxUMLClass[enumclass]{AICategory}
	{
		CHAT,
		CODEGENERATOR,
		POSTGENERATOR
	}

	\pxUMLClass{Employee}
	{
		- role : Role,
		- email : String
	}[+ monthlyConsumptionCost(month : DateTime) : real]

	\pxUMLClass{TokenConsumption}
	{
		- timeStamp : DateTime,
		- numTokens : UnlimitedNatural
	}

	\pxUMLClass{AIAgent}
	{
		- name : String,
		- description : String,
		- tokenPrice : Real,
		- aiCategory : AICategory
	}[+ process(prompt : String) : String, + calculateCost(numTokens : UnlimitedNatural):Real	]

	\pxUMLClass{Company}
	{
		- name : String,
		- address : String,
		- tel : String,
		- web : String
	}

	\pxUMLDiagram{
		\drawclass*[text width=15em]{Role}[2,0]
		\drawclass*[text width=14em]{AICategory}[16,0]
		\drawclass[text width=16em]{TokenConsumption}[9,-5]
		\drawclass[text width=16em]{AIAgent}[18,-5]
		\drawclass[text width=16em]{Employee}[0,-5]
		\drawclass[text width=16em]{Company}[0,-10]
		\association{Employee}{0..*}{ }{Company}{}{1}
		\association{Employee}{1}{ }{TokenConsumption}{}{0..*}
		\association{TokenConsumption}{0..*}{}{AIAgent}{}{1}
	}[1.1\linewidth]
	\section{Proves bàsiques de pxGDX}

	Configuració bàsica de la portada i el layout mitjançant \texttt{\textbackslash PECTeXconfig}:




	\sectionuoc{Primera secció UOC}
	Text d'exemple dins de la secció amb estil UOC.

	\subsectionuoc{Subsecció UOC}
	\subsubsectionuoc{Subsubsecció UOC}
	Més text amb l'estil de subsecció UOC.

	Una regla horitzontal fina:\\
	\pxHRule

	Una regla horitzontal gruixuda centrada a la pàgina:\\
	\pxPageRule*

	\pxIMG[.2\linewidth]{img/testimage}[primarycolor]{Patatas patatodos}[TestPicture]

	\pxIMG[.2\linewidth]{testimage}{Patatas patatodos}[TestPicture2]
	\pxIMG[.6\linewidth]{wrongimage}[primaryColor]{Patatas patatodos}[TestPicture2]

	\pxIMGpair{img/testimage}{img/testimage}[primarycolor]{Patatas patatodos}[TestPicture3]
	\pxIMGList{testimage}{ITERATION 1,ITERATION 2,ITERATION 3,ITERATION 4}

	\pxGrow{Patata}[\linewidth]
\end{document}